\section*{\texorpdfstring{\S\CurrentExerciseGroup}{Section \CurrentExerciseGroup}}
\phantomsection
\addcontentsline{toc}{section}{\texorpdfstring{Exercises for \S\CurrentExerciseGroup}{Exercises for Section \CurrentExerciseGroup}}

\begin{enumerate}
\def\labelenumi{\arabic{enumi})}
\item
  Let X be a locally compact subspace of a locally compact space T. Show that for every measure \(\lambda\) on T, the support of the measure \(\lambda_{X}\) is the trace on X of the support of the measure \(\varphi_{X} \cdot \lambda\) .
\item
  Let \(\mathrm{T}\) be a locally compact space, \(\mathrm{X}\) a locally compact subspace of \(\mathrm{T}\) , and \(\mu\) a positive measure on \(\mathrm{T}\) ; set \(\nu = \varphi_{\mathrm{X}} \cdot \mu = j(\mu_{\mathrm{X}})\) ( \(j\) the canonical injection of \(\mathrm{X}\) into \(\mathrm{T}\) ).
\end{enumerate}

\begin{enumerate}
\def\labelenumi{\alph{enumi})}
\item
  Show that every \(\mu_{\mathrm{X}}\) -negligible subset of \(\mathbf{X}\) is \(\nu\) -negligible.
\item
  For every numerical function \(g \geqslant 0\) defined on \(X\) , show that \(\int^{*} g \, d\mu_{X} = \int_{X}^{*} g \, d\nu\) (use \(a\) ), as well as formula (7) of §3).
\item
  Let \(\mathbf{g}\) be a function defined on \(X\) , with values in \(\overline{\mathbf{R}}\) or in a Banach space, \(\mathbf{g}'\) the extension of \(\mathbf{g}\) to \(T\) equal to 0 on \(T - X\) . For \(\mathbf{g}\) to be \(\mu_{X}\) -integrable, it is necessary and sufficient that \(\mathbf{g}'\) be \(\nu\) -integrable, in which case \(\int \mathbf{g} d\mu_X = \int_X \mathbf{g} d\nu\) (use \(b\) )).
\end{enumerate}

\begin{enumerate}
\def\labelenumi{\arabic{enumi})}
\setcounter{enumi}{2}
\tightlist
\item
  The notations being the same as in Exer. 2, assume that \(\mathbf{X}\) is an open subset of \(\mathbf{T}\) .
\end{enumerate}

\begin{enumerate}
\def\labelenumi{\alph{enumi})}
\item
  For every numerical function \(g \geqslant 0\) defined on \(X\) , show that \(\int^{*} g d\mu_{X} = \int_{X}^{*} g d\mu\) (use Exer. 2, Prop. 3 of §5 and Prop. 4 of §1).
\item
  Let g be a function defined on X, with values in \(\overline{R}\) or in a Banach space, \(g'\) the extension of g to T equal to 0 on T-X. For g to be \(\mu_{X}\) -integrable, it is necessary and sufficient that \(g'\) be \(\mu\) -integrable, in which case \(\int g d\mu_{X} = \int_{X} g d\mu\) (use a)).
\end{enumerate}

\begin{enumerate}
\def\labelenumi{\arabic{enumi})}
\setcounter{enumi}{3}
\item
  The notations being those of Exer. 2, show that if X is closed in T then, for every numerical function \(f \geqslant 0\) defined on T, \(\int^{*}(f \circ j) d\mu_{X} = \int^{*} f d\nu\) (observe that the mapping \(j\) is proper, and use Prop. 2 of §4). Under the same conditions, if \(\mathbf{f}\) is a mapping of T into \(\overline{\mathbf{R}}\) or into a Banach space, for \(\mathbf{f} \circ j\) to be \(\mu_{X}\) -integrable it is necessary and sufficient that \(\mathbf{f}\) be \(\nu\) -integrable, in which case \(\int(\mathbf{f} \circ j)\mu_{X} = \int \mathbf{f} d\nu\) (§4, Th. 2).
\item
  Let \(\mathbf{T}\) and \(\mu\) be the locally compact space and the measure defined in Exer. 5 of Ch. IV, §1, and let \(\mathbf{D}\) be the (closed) set of points of \(\mathbf{T}\) of the form \((0,y)\) .
\end{enumerate}

\begin{enumerate}
\def\labelenumi{\alph{enumi})}
\item
  Show that the measure \(\mu_{\mathrm{D}}\) induced on D by \(\mu\) is zero; however, \(\int^{*}\varphi_{\mathrm{D}}d\mu = +\infty\) (cf.~Exer. 3 a)).
\item
  Let \(\mathrm{X} = \mathrm{T} - \mathrm{D}\) , and let \(j\) be the canonical injection of \(\mathrm{X}\) into \(\mathrm{T}\) ; then \(j(\mu_{\mathrm{X}}) = \mu\) , but \(\int^{*} (\varphi_{\mathrm{D}} \circ j) d\mu_{\mathrm{X}} = 0\) , \(\int^{*} \varphi_{\mathrm{D}} d\mu = +\infty\) (cf.~Exer. 4).
\end{enumerate}

\begin{enumerate}
\def\labelenumi{\arabic{enumi})}
\setcounter{enumi}{5}
\item
  Let T be a locally compact space, X a locally compact subspace of T, and j the canonical injection of X into T. Let \(\lambda\) be a positive measure on X such that, for every compact set \(K \subset T\) , \(K \cap X\) is essentially \(\lambda\) -integrable. Show that the measure \(\nu = j(\lambda)\) is the smallest positive measure \(\rho\) on T such that \(\rho_{X} = \lambda\) , and that its support is the closure in T of the support of \(\lambda\) .
\item
  Let X, Y be two locally compact spaces, \(\pi: X \to Y\) a continuous mapping, T a locally compact subspace of Y, \(S = \overline{\pi}^{-1}(T)\) , \(\pi_{T}: S \to T\) the mapping that coincides with \(\pi\) in S, and \(\nu\) a positive measure on Y. Show that the following three conditions are equivalent:
\end{enumerate}

\begin{enumerate}
\def\labelenumi{\alph{enumi})}
\item
  If \(\nu_{\mathrm{T}}\) is the measure induced by \(\nu\) on \(\mathbf{T}\) , there exists a positive measure \(\mu_{\mathrm{T}}\) on \(\mathbf{S}\) such that \(\pi_{\mathrm{T}}(\mu_{\mathrm{T}}) = \nu_{\mathrm{T}}\) .
\item
  There exists a positive measure \(\lambda\) on \(X\) such that \(\pi (\lambda) = \nu_{\mathrm{T}}\) .
\item
  There exists a positive measure \(\lambda_{\mathrm{S}}\) on S such that \((\pi|\mathrm{S})(\lambda_{\mathrm{S}})=\nu_{\mathrm{T}}\) .
\end{enumerate}

\begin{enumerate}
\def\labelenumi{\arabic{enumi})}
\setcounter{enumi}{7}
\item
  Let X, Y be two locally compact spaces, \(\pi: X \to Y\) a continuous mapping, and \(\nu\) a positive measure on Y that is concentrated on \(\pi(X)\) . Assume that for every \(y \in \pi(X)\) , there exist a neighborhood V of y in Y and a positive measure \(\lambda_{\mathrm{V}}\) on \(\overline{\pi}^{-1}(V)\) such that the image of \(\lambda_{V}\) under the mapping \(\pi_{\mathrm{V}}: \overline{\pi}^{-1}(V) \to V\) that coincides with \(\pi\) is equal to the induced measure \(\nu_{V}\) . Show that under these conditions, there exists a positive measure \(\mu\) on X such that \(\pi(\mu) = \nu\) . (Consider the set G of pairs \((G, \lambda)\) , where G is open in Y and \(\lambda\) is a positive measure on X such that \(\pi(\lambda) = \nu_{\mathrm{G}}\) . Order G by the relation \(\langle G_{1} \subset G_{2} \text{ and } \lambda_{1} \leqslant \lambda_{2} \rangle\) between \((G_{1}, \lambda_{1})\) and \((G_{2}, \lambda_{2})\) . First show that the ordered set G is inductive, then complete the argument with the help of Zorn's lemma and Exer. 7.)
\item
  Let X, Y be two locally compact spaces, \(\pi: X \to Y\) a continuous mapping. The mapping \(\pi\) is said to be conservative if, for every positive measure \(\nu\) on Y that is concentrated on \(\pi(X)\) , there exists a positive measure \(\mu\) on X such that \(\nu = \pi(\mu)\) .
\end{enumerate}

\begin{enumerate}
\def\labelenumi{\alph{enumi})}
\item
  Let X,Y,Z be three locally compact spaces, \(\pi:X\to Y\) and \(\pi^{\prime}:Y\to Z\) two continuous mappings; set \(\pi^{\prime\prime}=\pi^{\prime}\circ\pi\) and assume that the mapping \(\pi\) is surjective. Show that if \(\pi\) and \(\pi^{\prime}\) are conservative, then so is \(\pi^{\prime\prime}\) ; conversely, if \(\pi^{\prime\prime}\) is conservative then so is \(\pi^{\prime}\) .
\item
  Assume that for every \(y \in \pi(X)\) there exists an open neighborhood V of y in Y such that, if \(\pi_{\mathrm{V}} : \overline{\pi}^{-1}(\mathrm{V}) \to \mathrm{V}\) is the mapping that coincides with \(\pi\) , there exists a continuous section associated with \(\pi_{V}\) (GT, I, §3, No.~5). Show that \(\pi\) is conservative (make use of Exer. 8).
\item
  Let Y be the interval \([0,1]\) of R, X the set Y equipped with the discrete topology, and \(\pi:X\to Y\) the identity mapping. Show that the mapping \(\pi\) is not conservative.
\end{enumerate}

\begin{enumerate}
\def\labelenumi{\arabic{enumi})}
\setcounter{enumi}{9}
\item
  Let \(X, Y\) be two locally compact spaces. Show that every proper continuous mapping \(\pi\) of \(X\) into \(Y\) is conservative. (The hypothesis implies that the mapping \(f \mapsto f \circ \pi\) of \(\mathcal{K}(Y)\) into \(\mathcal{C}(X)\) has its image \(E\) contained in \(\mathcal{K}(X)\) ; then, for every positive measure \(\nu\) on \(Y\) , apply to the subspace \(E\) of \(\mathcal{K}(X)\) and to the positive linear form \(f \circ \pi \mapsto \nu(f)\) on \(E\) the Cor. of Prop. 1 of TVS, II, §3, No.~1.)
\item
  Let X, Y be two locally compact spaces, \(\pi: X \to Y\) a continuous mapping, and M a subset of Y such that \(\overline{\pi}^{-1}(M)\) is contained in a countable union of compact sets. Then, for every positive measure \(\nu\) on Y that is concentrated on M, there exists a positive measure \(\mu\) on X such that \(\pi(\mu) = \nu\) (use Exer. 10).
\end{enumerate}
% semantic-fragments: legacy-input-seam
\relax{}
\setcounter{exercisegroup}{8}
\edef\CurrentExerciseGroup{\arabic{exercisegroup}}
